\section{La trayectoria de Liu Xianming (刘先明): de CV/ML a Robotaxi y después a un fabricante de automóviles}

Esta sección examina primero por qué Liu Xianming llegó a XPeng (小鹏). Su trayectoria incluye un doctorado en UIUC, Facebook Connectivity Lab, Horizon Robotics (地平线) en Norteamérica, Cruise Robotaxi y, por último, XPeng. Lo que da continuidad a esta línea no es un «historial de cambios de empleo», sino una elección tecnológica mission driven: de usar imágenes satelitales para atender desastres y estimar la distribución de la población mundial, a la conducción autónoma para reducir la fatiga de los conductores y los riesgos de tránsito, y después a construir, dentro de un fabricante de automóviles, un ciclo cerrado de datos controlable.

\begin{figure}[H]
\centering
\includegraphics[width=0.96\textwidth]{liuxianming-trajectory.png}
\caption{La trayectoria de Physical AI de Liu Xianming: CV/ML, Facebook, Cruise, XPeng y Physical AI. Diagrama conceptual de elaboración propia, a partir de la conversación de 00:02:16--00:19:00.}
\end{figure}

\begin{knowledgebox}{Lectura de la figura: el hilo conductor de esta trayectoria son los problemas del mundo real}
De las imágenes satelitales para atender desastres al Robotaxi, y después a la flota de vehículos de producción en serie de XPeng, Liu Xianming se ha ocupado siempre de «cómo entra la IA en el mundo real». Esto explica por qué no ve la conducción autónoma solo como una función de las empresas automotrices, sino como parte de Physical AI.
\end{knowledgebox}

\subsection{Dos lecciones de Cruise}

Dos lecciones que Liu Xianming aprendió en Cruise influyeron en su rumbo posterior. La primera es la simplificación extrema y la Infra a gran escala. Para que un Robotaxi pase de demo a producto se necesitan Data Infrastructure, Training Infrastructure, una cadena de análisis de problemas y un sistema de publicación de software. La segunda es la Continuous Learning Machine, es decir, resolver problemas de forma continua mediante la iteración con datos. Lo de «tumbarse en Hawái a esperar que caigan monedas de oro» es solo una broma; lo que hay detrás es que el ciclo cerrado de datos permite que el sistema mejore de manera continua.

\begin{importantbox}{Cómo se traslada la experiencia de Cruise}
Al pasar del Robotaxi a un fabricante de automóviles, lo decisivo no es cambiar de escenario, sino trasladar «el retorno de datos + Infra + la iteración rápida» a una flota más grande y a una cadena de datos más controlable.
\end{importantbox}

\subsection{Por qué ir a un fabricante de automóviles}

Si el futuro de la conducción autónoma y de Physical AI depende de datos a gran escala, de cadenas de datos controlables y de retroalimentación rápida, los fabricantes de automóviles tienen una ventaja natural. Las empresas de Robotaxi controlan su flota, pero su escala es limitada; los fabricantes tienen vehículos de producción en serie, usuarios reales, vialidades de distintas ciudades y una cadena completa de datos en el vehículo. Liu Xianming eligió XPeng porque él y He Xiaopeng (何小鹏) se hacían la misma pregunta: ¿cómo puede la ruta tecnológica de la siguiente generación dejar muy atrás a los competidores actuales?

\begin{knowledgebox}{Ventajas distintas de un fabricante de automóviles frente a una empresa de Robotaxi}
\begin{tabularx}{\textwidth}{p{0.22\textwidth}p{0.34\textwidth}X}
\toprule
Dimensión & Empresa de Robotaxi & Fabricante de automóviles \\
\midrule
Control de la flota & Controla la flota en operación; datos de alta calidad, pero la escala crece más despacio & Gran escala de vehículos de producción en serie; cobertura más amplia de ciudades y escenarios. \\
Objetivo comercial & Ofrece directamente servicios de movilidad & Venta de vehículos, suscripción de conducción inteligente, Robotaxi y experiencia de marca en paralelo. \\
Cadena de datos & Ciclo cerrado operativo sólido & El retorno de datos de la flota de usuarios y los problemas de la producción en serie son más complejos. \\
Restricciones de riesgo & Zona de servicio y alcance operativo controlables & Se dirige al usuario común; la responsabilidad de seguridad y calidad es mayor. \\
\bottomrule
\end{tabularx}
\end{knowledgebox}

\subsection{La presión que imponen los problemas del mundo real}

Liu Xianming comparó la presión laboral en Facebook, Cruise y XPeng. En Facebook, los problemas de investigación eran sobre todo de tecnología y recursos; en Cruise, los vehículos estaban en la calle, y los accidentes, las fallas y la seguridad eran presiones reales; en XPeng el ritmo es aún más rápido, porque un fabricante de producción en serie debe atender al mismo tiempo la IA, el producto, el negocio, la calidad, el hardware y la responsabilidad ante los usuarios. Esta presión explica por qué Physical AI no es research puro, sino que exige tomar decisiones de compromiso de forma constante dentro de sistemas reales.

\begin{warningbox}{De dónde viene la presión sobre la IA en el mundo real}
Cuando un modelo del mundo digital se equivoca, normalmente se puede revertir, reintentar u ocultar la falla; cuando un vehículo en la calle se equivoca, están en juego la seguridad, la responsabilidad, la marca y la regulación. Por eso la iteración de la IA física tiene que ser más cautelosa que la del software puro y depender más del ciclo cerrado de ingeniería.
\end{warningbox}

\subsection{De San Francisco a Cantón: el escenario cambia la estructura del problema}

Liu Xianming mencionó que, al trabajar en conducción autónoma en San Francisco, los corner case como personas en situación de calle, basura, mascotas y aves eran muy frecuentes, lo que hacía sentir que el problema era muy difícil; en cambio, en Cantón, algunos problemas que antes costaba resolver no tienen la misma magnitud. Esta observación es importante: la conducción autónoma no es un ejercicio abstracto de algoritmos; la estructura vial de la ciudad, los participantes del tránsito, la infraestructura y el comportamiento de la gente modifican la dificultad del problema.

\begin{importantbox}{El escenario no es un telón de fondo, sino parte del modelo}
El mismo algoritmo enfrenta corner case distintos en ciudades distintas. La elección del escenario afecta la distribución de los datos, el entrenamiento del modelo, el ritmo de puesta en producción y la experiencia del usuario; por eso, ni en el Robotaxi ni en la conducción inteligente de producción en serie se puede hablar solo del modelo sin hablar de la ciudad y del entorno operativo.
\end{importantbox}

\subsection{Resumen de la sección}

La trayectoria de Liu Xianming muestra que este cambio de liderazgo en XPeng no es solo un nombramiento organizacional, sino una nueva combinación de la experiencia en conducción autónoma, el ciclo cerrado de datos del Robotaxi y los recursos de producción en serie de un fabricante de automóviles.
